\chapter{The Concurrency Contract}
\label{ch:contract}

\standfirst{This chapter is normative. It states the one guarantee this system
takes away from Smalltalk-80, what replaces it, and what is still promised.
Read it before writing any Smalltalk that runs on more than one worker.}

\section{The short version}

\begin{stwarn}[The break]
\textbf{Smalltalk-80 guarantees that a process is never preempted by another
process of the same priority. Smalltalk-2026 does not.} Two processes at the
same priority may execute simultaneously on different CPUs. That is the entire
point of the system.
\end{stwarn}

\section{Why it matters more than it sounds}

The Blue Book scheduler contract reads: \emph{the scheduler is cooperative
between processes of the same priority, and preemptive between processes of
different priorities.}

In a green-threaded Smalltalk that is not merely a scheduling policy. It is an
\textbf{implicit mutual-exclusion primitive}, and forty years of Smalltalk code
leans on it everywhere. Every unsynchronised read-modify-write of an instance
variable, a class variable or a shared collection is correct \emph{only}
because no other process at the same priority can run between two bytecodes.

\begin{st}
"correct in every other Smalltalk; a lost update here"
count := count + 1
\end{st}

Code that raises its priority to get an atomic section is using the same trick
deliberately. Under true parallelism that trick stops working silently:
raising your priority no longer stops anybody, because they are on another
core.

This is why the class library has to be audited rather than inherited.

\section{What is guaranteed}

\begin{enumerate}
\item \textbf{Bytecode-level memory safety.} The object memory never tears. A
  field read yields some value previously written to that field, never a
  mixture. Object headers, the object table and the class of an object are
  always consistent.
\item \textbf{\ct{become:} is atomic.} Identity mutation is a single atomic
  swap of two object-table entries. No thread ever observes a half-completed
  \ct{become:}. This is the payoff for keeping an object table.
\item \textbf{\ct{Semaphore>>signal} and \ct{wait} are atomic} with respect to
  each other and to process suspension and resumption.
\item \textbf{Garbage collection is invisible.} Collection happens at
  safepoints, which the interpreter polls at message sends and backward jumps.
  No Smalltalk code observes a partially collected heap.
\item \textbf{Method lookup is coherent.} A class recompiled on one worker
  becomes visible to all workers; sends in flight complete against a
  consistent method dictionary.
\end{enumerate}

\section{What is not guaranteed}

\begin{enumerate}
\item \textbf{Same-priority atomicity.} Gone.
\item \textbf{Priority as mutual exclusion.} Priorities are scheduling hints.
  A higher-priority process is \emph{preferred}, not exclusive.
\item \textbf{Thread safety of the base collections.}
  \ct{OrderedCollection}, \ct{Dictionary}, \ct{Set} and \ct{Bag} are
  \textbf{not} synchronised. This follows Java's post-\ct{Vector} lesson:
  paying for a lock on every access to serve the rare shared case is the wrong
  default. Use the shared variants, or guard them yourself.
\item \textbf{\ct{Transcript} interleaving.} Output from concurrent processes
  may interleave unless you hold the Transcript's monitor across an entry.
\end{enumerate}

\section{What replaces it}

\begin{center}
\small
\begin{tabular}{@{}lp{0.60\textwidth}@{}}
\toprule
\ctp{Mutex} & non-reentrant mutual exclusion. \ctp{aMutex critical: [ ... ]}\\
\ctp{Monitor} & reentrant lock with condition variables:
  \ctp{waitUntil:}, \ctp{signalAll}\\
\ctp{SharedQueue} & multi-producer, multi-consumer, blocking\\
\ctp{Promise} & a value another process will supply. \ctp{aPromise value}\\
\ctp{Processor>>forkParallel:} & fork onto the worker pool and get a Promise\\
\bottomrule
\end{tabular}
\end{center}

\ct{Semaphore} remains, with Blue Book semantics, now genuinely atomic.

\section{The thread map}

\begin{plain}
Thread 0        Dedicated SDL pump.  Owns the window, renderer and texture.
                Never executes Smalltalk bytecodes.

Threads 1..N    Smalltalk workers, N ~ CPU count.  Run bytecodes, allocate
                from thread-local buffers, poll safepoints.  Never call SDL.

Delay timer     One request outstanding.  Sleeps on a condvar until the
                deadline, then posts to the async signal queue.  Never
                executes Smalltalk.
\end{plain}

Thread 0 is not a worker by design: if it were, it could be parked in a GC
safepoint at the moment the window server needs a response, deadlocking the
compositor. SDL3's own rules force the same shape---\ct|SDL_PumpEvents| and
friends are main-thread-only, and on macOS "main thread" means the thread that
entered \ct{main()}.

Smalltalk \ct{Process} objects remain green and cheap, multiplexed M:N over
the worker pool. \ct{Processor activeProcess} is per-worker state.

\begin{stnote}[What runs on the pool today]
That map is the design, and it is what the parallel test suites and
\ct{make bench} actually run---up to 31 native threads, ThreadSanitizer-clean,
as a gate rather than a report. It is \emph{not} yet what
\ct{./st80 -run image} does: the windowed desktop runs the interpreter on one
thread, so \ct{Processor workerCount} answers 1 there and in a
\ct{-bootstrap -eval} run.

That does not make this chapter advisory. The library your code will run
against is being audited to this contract now, and code written to it today
keeps working when the desktop moves onto the pool. Code written against
same-priority atomicity does not.
\end{stnote}

\section{The taxonomy to think in}

From Pallas and Ungar's \emph{Multiprocessor Smalltalk} (PLDI~'88), and this
system adopts it. For each piece of shared state, choose one:

\begin{center}
\small
\begin{tabular}{@{}lp{0.62\textwidth}@{}}
\toprule
\textbf{serialize} & lock it. Correct, and the answer of last resort---if
  everything lands here the system will not scale\\
\textbf{replicate} & give each worker its own. The \ctp{Transcript} buffer is
  this\\
\textbf{reorganize} & remove the sharing. Method dictionaries becoming
  immutable once published is this\\
\midrule
\textbf{immutable} & written once at build, never after\\
\textbf{thread-confined} & only ever touched by one worker.
  \ctp{ScheduledControllers} is this\\
\bottomrule
\end{tabular}
\end{center}

The last two are what keep the audit finite.

\section{Rules that follow}

\begin{description}
\item[Never lazily initialise a class variable.]
  \ct|^Default ifNil: [Default := Foo new]| is check-then-act: two workers
  each build one and one of them is thrown away, or worse, both are used.
  Initialise eagerly in a class-side \ct{initialize}, which the loader runs at
  build time.
\item[Do not use priority as a lock.] Raising your priority stops nobody on
  another core.
\item[Guard every shared collection, or do not share it.] Prefer not sharing:
  give each worker its own and merge at the end.
\item[Release locks in an \ct{ensure:}.] A critical section that exits by
  exception or non-local return must still release. \ct{Mutex} does this for
  you, which is why it could not be written before exceptions existed.
\item[Never hold a lock across a safepoint poll.] Sends and backward jumps
  poll. A lock held across one can deadlock the collector.
\end{description}

\section{One worker or many}

The suites and the benchmark take a worker count; \ct{make bench} sweeps 1, 2,
4, 8 and up. Running the same work at one worker and at many is the sharpest
debugging step there is: if a fault goes away at one worker, it is a race, and
this chapter is the list of places to look.

\begin{stnote}
A \ct{-bootstrap -eval} run has one worker. Concurrency code written against
this contract still works there---it is the same scheduler---but nothing will
actually run in parallel, so a race will not show up. Test with \ct{-run} and
several workers, or with the unit suites, which run at up to 31 threads under
ThreadSanitizer.
\end{stnote}
